% 第6章 强化学习
\section{强化学习：执行、组合与做市}\label{sec:rl}

前面两族都在做预测：先猜未来收益，再把预测变成仓位。强化学习（RL）换了个思路，直接学“在什么状态下该做什么动作”，把预测和决策合成一个目标。这个转向对交易尤其自然，因为交易本来就是序贯决策——今天的仓位决定明天的可选空间，执行成本随下单节奏变化，而目标（终值效用）不是单期收益的简单加总。

\subsection{问题设定}

把交易写成马尔可夫决策过程 $(\mathcal{S},\mathcal{A},\mathcal{P},\mathcal{R},\gamma)$：状态 $s_t$ 含持仓、行情、特征，动作 $a_t$ 是下单/调仓/报价，奖励 $r_t$ 是区间收益扣除成本，目标是最大化期望折扣回报
\begin{equation}
  J(\theta) = \mathbb{E}_{\pi_\theta}\Big[\sum_{t=0}^{T}\gamma^t\, r_t\Big].
\end{equation}
形式上干净，但金融场景下这条式子埋着三个坑：$\mathcal{P}$（环境转移）未知且不可反复采样，$\mathcal{R}$ 里成本项一旦设错就被“利用”，而 $\gamma$ 与 $T$ 的设定会改变策略是短视还是长视。\textbf{RL 的危险不在于学不会，而在于它会精确地学会你设定的那个目标——包括目标里的错误。}

\subsection{三个主要落地方向}

\textbf{（一）最优执行}。这是 RL 最早、也最成功的应用。Nevmyvaka、Feng 与 Kearns（2006）\cite{nevmyvaka2006} 用 RL 学习在限价订单簿中的下单时机，是对 Almgren--Chriss 解析解的第一次数据驱动挑战。这一方向的优势在于：环境相对稳定、目标清晰（最小化执行成本）、且有成熟的基线（VWAP/TWAP 与 Almgren--Chriss）可对照。近年一批工作把它做成深度 RL，并在含真实成本与市场冲击的模拟器里评测。Cartea、Jaimungal 与 Penalva（2015）\cite{cartea2015} 的系统论述为这一方向提供了理论底座。

\textbf{（二）组合管理}。Jiang、Xu 与 Liang（2017）\cite{jiang2017} 提出用深度 RL 做组合管理，网络直接输出各资产的权重，绕开了“先预测再优化”的两段式。Deng 等（2017）\cite{deng2017} 提出直接 RL 做金融信号表示与交易，把特征学习与决策学习连成一段。这一类工作后来被开源框架 FinRL（Liu 等\cite{liu2021finrl}）系统化，FinRL 提供了环境、数据管线与多种 RL 基线，是这一族可复现性最好的资产之一。但需要指出：\textbf{FinRL 的价值主要在于“把流程标准化”，而不在于“证明 RL 能赚钱”}——用它跑出来的高收益，很大一部分取决于环境与成本设定，而非算法本身。

\textbf{（三）做市}。做市是 RL 的天然场：报价是连续动作，目标是赚取买卖价差同时管理库存风险。Gu\'eant 与 Manziuk（2019）\cite{guant2019} 用深度 RL 做做市，Spooner 等（2018）\cite{spooner2018} 用 RL 学习做市报价，均报告在模拟环境中优于传统报价基线。做市方向的可信度相对高，因为目标（库存调整后的价差收益）与模拟环境（一笔笔订单流）能做得比较接近现实。

\subsection{落地要素}

\begin{itemize}
  \item \textbf{模拟器就是一切}。RL 学出的策略的好坏，几乎完全取决于训练环境与真实环境的接近程度。用历史数据回放（backtest-as-environment）训练，等于让策略去拟合“如果我的下单不影响市场”这个假世界。
  \item \textbf{成本与冲击必须内生化}。若把成本当成固定费率，策略会学会疯狂交易；只有让冲击随参与率上升，策略才会学会“慢下来”。
  \item \textbf{奖励设计要抗 Hack}。只给收益奖励，策略会学会承担尾部风险换取高均值；只给夏普，又会在样本内择时。
  \item \textbf{样本效率}。金融历史只有一份，RL 需要在有限交互下学到稳定策略，这与游戏/机器人动辄上亿次交互的场景根本不同。
\end{itemize}

\subsection{复现状态与失效边界}

FinRL\cite{liu2021finrl} 与若干做市/执行仓库提供了代码与数据处理，复现状态为\textbf{部分复现}：流程可跑通，但原文报告的收益在换数据、换成本口径后常不能保持。失效边界主要有三条。第一，\textbf{环境错配}——训练用的历史回放让策略“假设自己不影响市场”，实盘一放大参与率就失效。第二，\textbf{过拟合到制度}——RL 在单一制度下学到的策略在制度切换后崩溃，且因为它可以在大量参数组合里搜索，多重检验问题比监督学习更严重。第三，\textbf{奖励与真实目标的偏离}——奖励里的任何近似（成本、风险、折扣）都可能被策略放大成致命的行为。

本章要记住的一句话是：\textbf{强化学习在交易上最大的诱惑，是它把“预测--决策”一步到位；最大的风险，也是它把决策直接学进了你对世界的假设里——假设错了，它比你更彻底地错。}
